\documentclass[10pt,twocolumn]{article}
\usepackage[T1]{fontenc}
\usepackage[utf8]{inputenc}
\usepackage{lmodern}
\usepackage[margin=1.8cm,top=2.4cm,bottom=2cm,columnsep=0.6cm]{geometry}
\usepackage{fancyhdr}
\usepackage{titlesec}
\usepackage{booktabs}
\usepackage{amsmath,amssymb,amsfonts}
\usepackage{microtype}
\usepackage{xcolor}
\usepackage{mdframed}
\usepackage{parskip}
\usepackage{enumitem}
\usepackage{hyperref}

\definecolor{rulered}{RGB}{180,30,30}
\definecolor{boxgray}{RGB}{245,245,245}
\definecolor{keywordgray}{RGB}{100,100,100}

\pagestyle{fancy}
\fancyhf{}
\fancyhead[L]{\textsc{\small Claude Mag}}
\fancyhead[C]{\small\textit{Machine Learning Research Digest}}
\fancyhead[R]{\small 2026-09-23}
\fancyfoot[C]{\small\thepage}
\renewcommand{\headrulewidth}{0.8pt}

\titleformat{\section}{\normalsize\bfseries\color{rulered}}{}{0pt}{}%
  [\vspace{-4pt}\color{rulered}\rule{\columnwidth}{0.4pt}]
\titlespacing{\section}{0pt}{8pt}{4pt}

\hypersetup{colorlinks=true,urlcolor=rulered,linkcolor=black}

\begin{document}
\setlength{\parindent}{0pt}
\setlength{\parskip}{4pt}

{\small\color{keywordgray}\textbf{Keywords:}
  RLVR \textbullet{}
  entropy reward \textbullet{}
  LLM reasoning efficiency \textbullet{}
  chain-of-thought}\par\vspace{4pt}

{\LARGE\bfseries\raggedright
  ERR{+}: Sequential Entropy Resolution}\par\vspace{4pt}

{\large\itshape\raggedright
  Rewarding Token-Level Entropy Drops Teaches Decisive Reasoning and
  Cuts Chain-of-Thought Length on AIME, AMC, and STEM Benchmarks}\par\vspace{6pt}

{\small\textbf{By} Juxin Chen et al.}\par\vspace{2pt}

{\small\color{keywordgray}arXiv:2608.28771 \textbullet{} Findings of EMNLP 2026}
\par\vspace{2pt}\noindent{\color{rulered}\rule{\columnwidth}{0.6pt}}\vspace{6pt}

Correct chain-of-thought traces collapse token-level entropy more
frequently and more sharply than incorrect ones. ERR+ rewards these drops
in Phase~1 and compresses trace length in Phase~2, improving accuracy by
2--5 points and reducing trace length by 15--25\% on five math and STEM
benchmarks.

\begin{mdframed}[backgroundcolor=boxgray,linecolor=rulered,linewidth=1pt,
    innerleftmargin=6pt,innerrightmargin=6pt,
    innertopmargin=6pt,innerbottommargin=6pt]
\textbf{\small AT A GLANCE}\par\vspace{4pt}
\begin{description}[leftmargin=0pt,labelsep=4pt,itemsep=2pt,parsep=0pt,
    font=\small\bfseries]
  \item[Problem:] RLVR-trained large reasoning models produce accurate but
    verbose traces; existing length penalties cut indiscriminately and
    degrade accuracy.
  \item[Why it matters:] Long traces waste inference compute; for deployed
    LRMs on math and STEM tasks, shorter decisive reasoning directly reduces
    serving cost.
  \item[Method:] Two-phase RLVR: Phase~1 rewards log-normalized cumulative
    entropy drops (gated on correctness); Phase~2 compresses length
    relative to rollout-group peers via GRPO.
  \item[Result:] +2--5 pp accuracy, --15--25\% trace length on AIME,
    AMC, MATH-500, GPQA-Diamond, STEM multi-step.
  \item[Evidence:] Ablations across 4 design axes confirm both phases
    are necessary; log normalization critical for preventing length gaming.
\end{description}
\end{mdframed}

\section{The Big Picture}

Reinforcement learning with verifiable rewards (RLVR) has produced
state-of-the-art reasoning models. These models generate extended
chain-of-thought (CoT) traces---sometimes spanning thousands of tokens---
before emitting a final answer. The length is partly necessary: complex
multi-step problems genuinely require extended computation. But a
substantial fraction of trace tokens represents redundant backtracking,
repeated reformulation, and exploratory text that does not advance the
reasoning.

The tension is delicate. Simply penalising length reduces redundancy but
also cuts useful deliberation, leading to accuracy regression. What is
needed is a way to distinguish \emph{decisive} reasoning steps---moments
when the model resolves uncertainty and commits to a path---from
\emph{indecisive} ones that wander or repeat.

ERR+ proposes that token-level entropy drops are precisely this signal.
When the model collapses its next-token distribution sharply, it has made
a committed inference. Correct reasoning chains do this more frequently
than incorrect ones, and rewarding this pattern produces models that are
simultaneously more accurate \emph{and} more concise.

\section{Why Existing Approaches Fall Short}

\textbf{Binary outcome RLVR} rewards all correct traces equally,
regardless of whether they are decisive or verbose. It provides no
gradient signal toward more efficient reasoning.

\textbf{Length penalties} (absolute or relative) reduce trace length but
penalise useful long reasoning steps alongside useless ones. The accuracy
degradation is well-documented.

\textbf{KL regularisation toward a reference model} prevents entropy
collapse globally but does not specifically encourage entropy drops
\emph{during} the thinking phase---it regulates the overall distribution,
not the dynamics within a trace.

\textbf{Min-token objectives} force early answer commitment, degrading
accuracy on hard problems that genuinely require extended reasoning.

ERR+ is the first method to use the temporal \emph{pattern} of entropy
changes within the thinking phase as a training signal, rather than
aggregate trace length or final-state entropy.

\section{Core Mechanism}

\textbf{The Entropy Drop Observation.} Let $H_t = -\sum_v p_\theta(v|x_{<t})
\log p_\theta(v|x_{<t})$ be the next-token entropy at position $t$ within
the thinking phase $\mathcal{T}_{\mathrm{think}}$. An entropy \emph{drop}
at position $t$ is $\Delta H_t = H_{t-1} - H_t > 0$.

Empirical analysis across multiple model families (7B--70B) reveals:
\begin{itemize}[itemsep=2pt,parsep=0pt]
  \item Correct reasoning traces have 30--50\% more entropy drops on
    average than incorrect ones on MATH and AIME tasks
  \item The magnitude of drops is also larger in correct traces
  \item This pattern holds even when controlling for total trace length
\end{itemize}

\textbf{Phase~1: Entropy Relief Reward (ERR).} The cumulative entropy
drop, log-normalized and gated on correctness, forms the ERR bonus:
\[
\mathrm{ERR}(\tau)
= \frac{\sum_{t\in\mathcal{T}_{\mathrm{think}}}
      \max(0,\,H_{t-1}-H_t)}
      {\log(|\tau|+1)}
  \cdot \mathbf{1}[a(\tau)=a^*]
\]

\textbf{Phase~2: Group-Relative Length Compression.} A GRPO-style penalty
rewards traces that are shorter than their rollout-group average:
\[
r_{\mathrm{len}}(\tau)
= -\gamma\cdot\frac{|\tau| - \bar{|\tau|}_{\mathrm{group}}}
                   {\bar{|\tau|}_{\mathrm{group}}}
\]

\section{Technical Formulation}

Token-level entropy:
\[
H_t = -\sum_{v\in\mathcal{V}}
  p_\theta(v\mid x_{<t})\,\log p_\theta(v\mid x_{<t})
\]

Log-normalized entropy relief:
\[
\mathrm{ERR}(\tau)
= \frac{1}{\log(|\tau|+1)}
  \sum_{t\in\mathcal{T}_{\mathrm{think}}}
  \max(0,\,H_{t-1}-H_t)
  \cdot \mathbf{1}[a(\tau)=a^*]
\]

Phase~1 total reward:
\[
r_1(\tau) = r_{\mathrm{base}}(\tau) + \beta\cdot\mathrm{ERR}(\tau)
\]

Phase~2 total reward:
\[
r_2(\tau) = r_{\mathrm{base}}(\tau)
  - \gamma\cdot\frac{|\tau|-\bar{|\tau|}_{\mathrm{group}}}
                    {\bar{|\tau|}_{\mathrm{group}}}
\]

Phase~1 policy gradient:
\[
\nabla_\theta\mathcal{L}_1
= -\mathbb{E}_\tau\!\left[r_1(\tau)\cdot
  \nabla_\theta\log p_\theta(\tau)\right]
\]

\section{Learning or Inference Procedure}

\begin{enumerate}[itemsep=2pt,parsep=0pt]
  \item Start from a CoT-fine-tuned base model.
  \item \textbf{Phase~1:} Run RLVR with $r_1$ for $N_1$ steps;
    monitor accuracy and mean CED (cumulative entropy drop).
  \item \textbf{Phase~2:} Continue from Phase~1 checkpoint with $r_2$
    for $N_2$ steps; monitor accuracy and trace length.
  \item Evaluate at fixed inference budget (greedy, no search).
\end{enumerate}

\section{What the Guarantee Says}

The log-normalization in ERR provides a formal argument against length
gaming: the ERR bonus grows only as $O(\log |\tau|)$ with length, while
the number of possible entropy drops grows as $O(|\tau|)$. A model that
increases length purely to accumulate entropy drops earns a
$O(|\tau| / \log|\tau|)$ bonus but pays increasing KL cost---net loss
for long random traces.

\section{Experimental Findings}

\begin{itemize}[itemsep=2pt,parsep=0pt]
  \item \textbf{Benchmarks:} AIME 2024, AMC 2023, MATH-500,
    GPQA-Diamond, custom STEM multi-step dataset
  \item Accuracy improves by 2--5 pp over RLVR-only baseline
  \item Trace length decreases by 15--25\% relative to Phase~1 baseline
  \item ERR+ outperforms length-penalty-only methods at matched length
  \item Mean CED increases by $\approx$40\% after Phase~1 training,
    confirming decisiveness is learned, not length
\end{itemize}

\section{Ablations and Interpretation}

\textbf{Phase~1 only:} Accuracy improves; trace length unchanged.
ERR alone drives decisiveness but not compression.

\textbf{Phase~2 only (no Phase~1):} Length decreases; accuracy
\emph{degrades} by 3--4 pp. Phase~1 first establishes the decisive
reasoning style that Phase~2 then compresses.

\textbf{No log-normalization:} Accuracy improves; trace length
\emph{increases}---the model games the reward by generating more drops
absolutely, confirming that normalization is essential.

\textbf{No correctness gating:} Entropy drops in incorrect traces are
also rewarded; reasoning quality degrades, confirming that decisive wrong
reasoning is worse than exploratory wrong reasoning.

\textbf{$\beta$ sweep:} Optimal $\beta\in[0.2, 0.5]$; large $\beta$
collapses entropy across all positions, hurting exploration.

\section*{Reference}

Juxin Chen et al.
\textbf{ERR{+}: Sequential Entropy Resolution for Efficient and
Decisive LLM Reasoning}.
arXiv:2608.28771. Findings of EMNLP 2026.\\
\url{https://arxiv.org/abs/2608.28771}

\end{document}
